\chapter{От микроуровня к макроописанию}
\label{ch:macro}

\section{Необходимость перехода}

Планковские $\caStepSpace$ и ячейка~$\PlanckCell$ определены через $\planckHbar$, $\constGrav$, $\speedC$,
но принципиально недоступны непосредственному измерению: слишком малы масштаб и объём.
Классическая физика часто берёт предел $c\to\infty$ или $\caStepSpace\to 0$;
на микроуровне~$M$ такой предел отвергнут (\cref{sec:no-continuum-carrier}).

Макроскопический мир~$T$ --- единственный слой прямого доступа наблюдателя.
Из~$M$ \emph{обязано} следовать~$T$: осреднение поля, дискретный анализ, статистика измерений.
Согласование модели с экспериментом проверяется на~$T$ (дисперсия, массы, Стандартная модель);
микроуровень подтверждается косвенно: если макроинварианты сходятся,
каркас не опровергнут.

\section{Электрон как составная структура}

Центральный вопрос перехода $M\to T$: из чего на микроуровне состоит
макроскопический электрон с известными массой~$\massElectron$ и зарядом~$\chargeE$?
Предел масштаба --- ячейка~$\PlanckCell$.
Проверка теории --- способность комбинаций ячеек и закона~$g$
воспроизвести электрон без новых подгоняемых параметров.

\section{Оператор осреднения \texorpdfstring{$\mathcal{B}$}{B}}
\label{sec:readout}

Вопрос не в «вкусе сглаживания'', а в том, какой ответ даёт прибор,
если микрошаг ограничен световым конусом (\cref{ax:A1}) и одна планковская ячейка
не является наблюдаемой (\cref{ch:foundations}).
Фиксируем четыре посылки прибора на~$\Tlayer$:
\begin{enumerate}[label=\textbf{C\arabic*:},leftmargin=*]
\item \label{cond:C0}
  прибор не разрешает одну ячейку~$\PlanckCell$ --- только поле на многих узлах;
\item \label{cond:C1}
  до статистики $|\Phi|^2$ осреднение линейно по~$z$ (макро-линейность, \cref{ax:A8});
\item \label{cond:C2}
  веса ядра неотрицательны и нормированы: $W\ge 0$, $\sum W=1$;
\item \label{cond:C3}
  за один такт~$\htick$ сигнал идёт только по светоподобным рёбрам длины~$\caStepSpace$;
  центральный узел за один такт в апертуру не попадает;
\item \label{cond:C4}
  вес узла в мягком поле пропорционален числу допустимых null-путей,
  дошедших до события прибора (дискретный Green конуса), а не произвольной «шапке''.
\end{enumerate}

\begin{definition}[Осреднение с биномиальным ядром]
\label{def:soft-readout}
На ортогональном срезе произведения осей
\[
  \Phi(X,t)=\bigl(\tfrac14[1,2,1]\bigr)^{\otimes 2R} * z(x,t).
\]
Эквивалентно одному проходу свёртки $3\times 3$ с ядром $[1,2,1;2,4,2;1,2,1]/16$.
Параметр $R\ge 1$ --- число проходов (ширина прибора), не планковский предел.
\end{definition}

\begin{proposition}[Вывод \texorpdfstring{$[1,2,1]/4$}{(1-2-1)} на цепи]
\label{prop:kernel-121}
Пусть на одномерной цепи~$\mathbb{Z}$ выполнены \ref{cond:C3}--\ref{cond:C4}
и глубина апертуры минимальна: ровно два такта (туда и обратно).
Тогда линейное осреднение (\ref{cond:C1}) с весами по числу путей даёт ядро
$W=\tfrac14[1,2,1]$ на соседях $\{-1,0,+1\}$.
\end{proposition}

\begin{proof}
\pusing{\ref{cond:C3}}
за один такт из $x$ достижимы только $x\pm\caStepSpace$; центр $x$ не входит в однотактовую апертуру.
\pusing{\ref{cond:C4} и минимальной глубины~$2$}
учитываем все пути длины~$2$ с возвратом в окрестность $x$.
Один light-like шаг вдоль оси --- равновероятный выбор $\pm 1$:
атом $w_1=\tfrac12[1,1]$.
Два независимых шага дают смещения $-2,0,+2$ с кратностями $1,2,1$.
После переноса весов на support $\{-1,0,+1\}$ относительно центра и нормировки (\ref{cond:C2})
получаем $w_1*w_1=\tfrac14[1,2,1]$.
\end{proof}

\begin{remark}[Ортогональный срез и \latGCKshort]
На product-срезе (независимые null-шаги по осям) тот же аргумент по каждой оси
и \ref{cond:C1} дают $\mathcal{B}$ из \cref{def:soft-readout}.
На полном носителе \latGCKshort\ соседство~$N_{12}$ не является кубом $(1{:}2{:}1)^{\otimes 3}$:
нативное ядро глубины~$2$ --- сумма по путям на графе~$N_{12}$
(в \cref{thm:t-cr}: $w(0)=1/12$, изотропный второй момент $\tfrac43\delta_{ij}$).
Сепарабельная формула в \cref{def:soft-readout} --- удобная запись на срезе;
спектральные оценки T-CR переносятся на ядро Грина по $N_{12}$ с той же шкалой~$\sigma_R$.
\end{remark}

\begin{remark}[Что не следует]
\begin{itemize}[nosep]
\item \textbf{Не} утверждается единственность среди всех положительных ядер с тем же support:
  выбран класс C0--C4 (Green путей), а не block-mean и не равномерный шар в решётке ---
  они дают ложные края и не совпадают со счётом null-путей.
\item При фиксированном $R\ge 1$ оператор $\mathcal{B}$ не инъективен
  (лемма «$\mathcal{B}$ необратима'', \cref{lem:236}): разные $z$ могут давать одно~$\Phi$.
\end{itemize}
\end{remark}

\section{Теорема T-CR: континуум как следствие осреднения}
\label{sec:t-cr}

\begin{theorem}[T-CR]
\label{thm:t-cr}
Континуум не постулат микроуровня~$M$.
При фиксированной $\ell_P$ и радиусе осреднения $R\ge 1$
любой прибор с ядром~$\mathcal{B}$ (условия C0--C4) видит поле,
неотличимое от непрерывного на шкалах $\lambda\gg\sigma_R\ell_P$,
с относительной ошибкой $O((\sigma_R\ell_P/\lambda)^4)$.
\end{theorem}

\begin{proof}
\puseref{осреднения с биномиальным ядром (C0--C4)}{def:soft-readout}
ядро $\mathcal{B}$ --- $R$-кратная свёртка с $W=\tfrac14[1,2,1]^{\otimes 2R}$ на ортогональном срезе.

\smallskip
\emph{Одномерная цепь.}
\pfrom{одного прохода $W=[1,2,1]/4$}
символ Фурье
$\widehat W(k)=\tfrac12+\tfrac12\cos k=\cos^2(k/2)$.
После $R$ проходов $\widehat W_R(k)=\cos^{2R}(k/2)$.
Один проход сдвигает узел на $-1,0,+1$ с вероятностями $1/4$, $1/2$, $1/4$,
поэтому дисперсия равна $1/4+1/4=1/2$.
$R$ проходов складывают дисперсии: $\sigma_R^2=R/2$.
Гауссово ядро с той же дисперсией имеет символ $G_R(k)=\exp(-\sigma_R^2 k^2/2)=\exp(-Rk^2/4)$.

\ptransform{разложения $\ln\cos u=-u^2/2-u^4/12+O(u^6)$ при $u=k/2$}
\[
  \ln\cos(k/2)=-\frac{k^2}{8}-\frac{k^4}{192}+O(k^6).
\]
Умножение на $2R$ даёт логарифм символа:
\[
  \ln\widehat W_R(k)=-\frac{Rk^2}{4}-\frac{Rk^4}{96}+O(Rk^6),
\]
и отношение к гауссову символу
\[
  \frac{\widehat W_R(k)}{G_R(k)}=\exp\Bigl(-\frac{Rk^4}{96}+O(Rk^6)\Bigr).
\]
При фиксированном~$R$ и $k\sigma_R\ll 1$ отклонение отношения от единицы есть $O(k^4)$,
то есть $O((k\sigma_R)^4)$.
Длина волны $\lambda=2\pi/k$ при $\ell_P=\caStepSpace$ переводит ту же оценку в $O((\sigma_R\ell_P/\lambda)^4)$.

\smallskip
\emph{Ультрафиолет и отсутствие $\caStepSpace\to 0$.}
При $k\sim 1$ (масштаб $\ell_P$) имеем $\widehat W_R(k)\le e^{-cR}$: осреднение подавляет
УФ-моды без предела $\caStepSpace\to 0$; континуум --- свойство осреднения на~$T$, не аксиома~$M$.

\smallskip
\emph{\latGCKshort\ в $(3{+}1)$.}
На \latGCK\ ядро глубины~$2$ по путям (\cref{def:soft-readout}, гл.~\ref{ch:macro})
даёт $w(0)=1/12$ и изотropный второй момент $M_{ij}=\tfrac43\delta_{ij}$
(полное перечисление двухшаговых путей на~$N_{12}$).
Separable накопление $W^{\otimes 2R}$ переносит одномерную оценку $O(k^4)$
на каждую ось product-среза; изотropия $M$ обеспечивает одну шкалу $\sigma_R$ во всех
направлениях осреднения.
\end{proof}

\section{Теорема T-CL: аналитический классический предел}
\label{sec:t-cl}

\begin{theorem}[T-CL]
\label{thm:t-cl}
На слое~$\Tlayer$ при фиксированных $\caStepSpace$, $\htick$ на~$\Mlayer$
«классика'' --- предел по масштабу наблюдателя, не по $\caStepSpace\to 0$.
Пусть $R$ --- радиус осреднения (\cref{def:soft-readout}),
$\sigma_R^2=R/2$ на product-срезе,
$\lambda$ --- длина волны наблюдаемой моды,
\[
  \varepsilon=\frac{\sigma_R\caStepSpace}{\lambda},\qquad
  \chi=\Bigl(\frac{\caStepSpace}{\lambda}\Bigr)^2.
\]
При $\varepsilon\ll 1$, $\chi\ll 1$ и посылках \cref{thm:t-cr} и гидропредела
(изотропное ядро глубины~$2$ на \latGCKshort):
\begin{enumerate}[label=(\roman*)]
\item \textbf{Спектр.} $|\widehat W_R(k)-G_R(k)|/G_R(k)=O(\varepsilon^4)$ при $k\sim 2\pi\caStepSpace/\lambda$.
\item \textbf{Лапласиан.} $(K^{*R}-I)\Phi=\tfrac{2R}{3}\caStepSpace^2\nabla^2\Phi+O(\varepsilon^4|\Phi|)$ в IR.
\item \textbf{Динамика.} Эффективное уравнение NLSE/GL с вязкостью $\nu$ и остатком $O(\varepsilon^4)$
  (разложение насыщающей фазы на медленном $\Phi$, \cref{sec:gl-free-energy}).
\item \textbf{Гидро.} Madelung-разложение даёт continuity+Euler/NS;
  квантовое давление Бома $Q_{\mathrm{Bohm}}=O(\chi)$ на гладких модах;
  при доминировании инерции и $\nu\nabla^2\mathbf v$ над $Q_{\mathrm{Bohm}}$ ---
  классическая гидродинамика.
\end{enumerate}
\end{theorem}

\begin{remark}[Про $\caStepSpace\to\infty$ и $t_P\to\infty$]
Увеличение $\caStepSpace$ при фиксированном метре эквивалентно $\varepsilon,\chi\to 0$
(больше ячеек на длину волны), а не деформации планковской ячейки на~$\Mlayer$.
Один лишь рост $t_P$ при фиксированном $\caStepSpace$ разрывает связку $\caSignalSpeed=\caStepSpace/\htick$
(\cref{ch:carrier}); временная «классика'' --- это много тактов~$g$ на один макро-такт прибора.
\end{remark}

\begin{definition}[Квантовое давление Бома на~$\Tlayer$]
\label{def:bohm-pressure}
Для плотности $\rho>0$ на макроскопическом слое положим $u=\sqrt{\rho}$.
В координатах прибора (единицы длины --- $\caStepSpace$ на~$\Mlayer$, см.~\cref{ch:carrier})
\[
  Q_{\mathrm{Bohm}}(x)\;:=\;-\frac{\planckHbar^2}{2m_{\mathrm{eff}}}\,
  \frac{\partial_x^2 u(x)}{u(x)},
\]
где $m_{\mathrm{eff}}$ --- эффективная масса пакета на~$\Tlayer$ из разложения насыщающей фазы
(\cref{sec:gl-free-energy}); $\partial_x$ --- производная по координате прибора.
В силовом уравнении Madelung на~$\Tlayer$ входит $-\partial_x Q_{\mathrm{Bohm}}$.
\end{definition}

\begin{lemma}[Оценка $Q_{\mathrm{Bohm}}$ на гладкой модуляции]
\label{lem:bohm-chi}
Пусть $\rho(x)=\rho_0\bigl(1+\eta\cos(2\pi x/\lambda)\bigr)$ с $0\le\eta<1$,
$\lambda\gg\caStepSpace$, $\chi=(\caStepSpace/\lambda)^2$.
Тогда при $\chi\ll 1$
\[
  \sup_x |Q_{\mathrm{Bohm}}(x)|
  \;\le\;
  \frac{\planckHbar^2}{m_{\mathrm{eff}}}\,
  \eta\,\chi\,\pi^2\,\bigl(1+O(\eta+\chi)\bigr).
\]
В частности $\sup|Q_{\mathrm{Bohm}}|=O(\eta\chi)=O(\chi)$ при фиксированной амплитуде $\eta$.
\end{lemma}

\begin{proof}
Безразмеризуем $s=x/\lambda$, $k=2\pi/\lambda$; положим $\rho/\rho_0=1+\eta\cos\theta$, $\theta=ks$,
$u(s)=\sqrt{1+\eta\cos\theta}$.
\phence $\partial_x=k\partial_\theta$, $\partial_x^2=k^2\partial_\theta^2$.

\smallskip
\emph{Разложение по $\eta$.}
При $|\eta|<1$ разложим $(1+\eta\cos\theta)^{1/2}$:
\[
  (1+\eta\cos\theta)^{1/2}
  =1+\tfrac{\eta}{2}\cos\theta-\tfrac{\eta^2}{8}\cos^2\theta+O(\eta^3),
\]
\pfollows
\[
  u=1+\frac{\eta}{2}\cos\theta+O(\eta^2),\qquad
  u_\theta=-\frac{\eta}{2}\sin\theta+O(\eta^2),\qquad
  u_{\theta\theta}=-\frac{\eta}{2}\cos\theta+O(\eta^2).
\]
Тогда
\[
  \frac{\partial_x^2 u}{u}
  =k^2\,\frac{u_{\theta\theta}}{u}
  =-\frac{k^2\eta}{2}\cos\theta+O(\eta^2 k^2),
\]
и
\[
  \Bigl|\frac{\partial_x^2 u}{u}\Bigr|
  \le\frac{k^2\eta}{2}+O(\eta^2 k^2).
\]
Подставляя $k=2\pi/\lambda$ и $\chi=(\caStepSpace/\lambda)^2$ при $\caStepSpace=1$ в решёточных единицах,
$k^2=4\pi^2\chi$, получаем
$\bigl|\partial_x^2 u/u\bigr|\le 2\pi^2\eta\chi+O(\eta^2\chi)$.
Умножение на $\planckHbar^2/(2m_{\mathrm{eff}})$ даёт заявленную границу с $O(\eta+\chi)$ в скобке.

\smallskip
\emph{Замечание о дискретной сетке.}
На конечной решётке $\partial_x$ в~\cref{def:bohm-pressure} можно понимать как центральную разность
шага $\caStepSpace$; при $\lambda\gg\caStepSpace$ замена $\partial_x^2\to\Delta^{(2)}_{\caStepSpace}$
меняет правую часть на ту же ведущую $O(\eta\chi)$ с иной константой
(ошибка дискретизации $O(\chi)$).
\end{proof}

\begin{proof}[Доказательство пункта~\emph{(iv)} теоремы~\ref{thm:t-cl}]
\pfrom{разложения Маделунга $\Phi=\sqrt\rho\,e^{iS/\planckHbar}$ на~$\Tlayer$}
получаем уравнение непрерывности для $\rho$ и уравнение импульса для $v=\partial_x S/m_{\mathrm{eff}}$
с силой $-\partial_x Q_{\mathrm{Bohm}}$ плюс градиенты давления от нелинейности насыщающей фазы и вязкости $\nu$
(модель, \S4.1.1-HL).

\pusing{моды $\rho=\rho_0(1+\eta\cos(2\pi x/\lambda))$ и леммы~\ref{lem:bohm-chi}}
\[
  \sup_x|\partial_x Q_{\mathrm{Bohm}}(x)|=O\!\left(\frac{\eta\chi}{\lambda}\right)=O\!\left(\frac{\chi}{\lambda}\right).
\]
Инерционный член $\rho(\partial_t v+v\partial_x v)$ и вязкий $\nu\rho\partial_x^2 v$ на том же масштабе
задаются макроскопической скоростью и $\nu$.
\phence при $\chi\to 0$ и фиксированных $\rho_0$, $|v|$ доминируют $\nu\partial_x^2 v$ и $\rho v\partial_x v$
над $-\partial_x Q_{\mathrm{Bohm}}$, и ведущим описанием пакета становится классическая гидродинамика (Навье--Стокс/Euler).
\end{proof}

\section{Свободная энергия Ginzburg--Landau}
\label{sec:gl-free-energy}

Полная микродинамика~$g$ на~$\Mlayer$ слишком детальна для макроскопического прибора.
Ситуация упрощается на шкалах $\lambda\gg\sigma_R\ell_P$ (\cref{thm:t-cr}): осреднение с биномиальным ядром (\cref{def:soft-readout})
оставляет медленное поле $\Phi=\mathcal{B}z$, которое меняется от узла к узлу слабо по сравнению
с фазовой структурой внутри~$\mathcal{B}$.

Различие структурированного и размазанного осреднённого поля удобно описывают параметром порядка.
Естественным выбором служит та же $\Phi$ (аналог $\psi$ в теории Гинзбурга--Ландау, см.~Landau--Lifshitz, т.~IX, \S45):
амплитуда и фаза на~$\Tlayer$.

\emph{Отправным пунктом} служит свободная энергия на объёме прибора~$V\subset\Tlayer$
как функционал от~$\Phi(\mathbf{x})$.
По общим соображениям теории Ландау о фазовых переходах второго рода
(там же, \S45; общая схема --- т.~V, \S146) плотность свободной энергии разлагается
по степеням малого $\Phi$ и его \emph{первых} производных.
Члены с $\nabla^2\Phi$, $\nabla^4\Phi$, \ldots{} отбрасываются: после~$\mathcal{B}$ поле достаточно
медленно в пространстве (аналог «фигурируют лишь первые производные'' в \S45).

Глобальная фаза $U(1)_{\mathrm{vac}}$ переводит $\Phi\mapsto\Phi\,e^{i\theta}$.
Физические величины на~$\Tlayer$ не должны зависеть от этой фазы с~$M$;
поэтому в разложении допустимы инварианты $|\Phi|^2$, $|\Phi|^4$, $|\nabla\Phi|^2$.

Нелинейность насыщающей фазы при $\rho_E\ll K_P$ (\cref{def:kp}, гл.~\ref{ch:evolution})
даёт ведущий член насыщения $\sim|\Phi|^4$; разложение вращателя по градиентам фазы
после~$\mathcal{B}$ --- член $\sim|\nabla\Phi|^2$.
Не повторяя полностью рассуждений гл.~\ref{ch:evolution}, запишем полную свободную энергию:
\begin{equation}
\label{eq:gl-free-energy}
  \mathcal{F}_{\mathrm{GL}}[\Phi]
  = \mathcal{F}_0
  + \int_V \Bigl(
      \alpha\,|\Phi|^2
      + \frac{\beta}{2}\,|\Phi|^4
      + \frac{\planckHbar^2}{2m}\,|\nabla\Phi|^2
    \Bigr)\,\mathrm{d}^3x.
\end{equation}
Здесь $\mathcal{F}_0$ --- свободная энергия однородного фона (при $\Phi\equiv 0$ на выбранном вакууме);
$\alpha$ --- квадратичный коэффициент (положение «дна'' потенциала на~$\Tlayer$);
$\beta>0$ --- коэффициент насыщения, в замыкании модели $\beta\equiv\omega_{\mathrm{macro}}$ из насыщающей фазы;
множитель $\planckHbar^2/(2m)$ связывает градиентную энергию с механической лестницей ---
по той же логике, что в \S4 т.~I выбор $L\propto mv^2$ фиксируется симметрией.

В однородном случае ($\Phi=\mathrm{const}$) минимизация~\eqref{eq:gl-free-energy} по $|\Phi|^2$
при $\alpha<0$ даёт ненулевое равновесное $|\Phi|^2=-\alpha/\beta$;
на фоне осциллирующего вакуума коэффициент~$\alpha$ фиксирован.

Варьируя~\eqref{eq:gl-free-energy} по $\Phi^*$ при фиксированных данных на $\partial V$:
\begin{equation}
\label{eq:gl-variation}
  \frac{\delta\mathcal{F}_{\mathrm{GL}}}{\delta\Phi^*}
  = \alpha\,\Phi
  - \frac{\planckHbar^2}{2m}\,\nabla^2\Phi
  + \beta\,|\Phi|^2\Phi.
\end{equation}
Условие $\delta\mathcal{F}_{\mathrm{GL}}=0$ --- статика медленного поля на~$\Tlayer$.

Для \emph{медленной} эволюции $\Phi(t)$ динамика на~$\Tlayer$ задаётся тем же функционалом:
\begin{equation}
\label{eq:gl-dynamics}
  i\planckHbar\,\frac{\partial\Phi}{\partial t}
  = \frac{\delta\mathcal{F}_{\mathrm{GL}}}{\delta\Phi^*}.
\end{equation}
Это не новый постулат на~$\Mlayer$ (там действует только~$g$ из \cref{sec:g-law}),
а макроскопическое уравнение, аналогичное переходу от $L$ к уравнениям Лагранжа
(т.~I, \S2) и от свободной энергии --- к уравнениям Гинзбурга--Ландау (т.~IX, \S45).

\section{Нелинейное уравнение Шрёдингера и гидродинамика}
\label{sec:nlse-hydro}

Подставляя~\eqref{eq:gl-variation} в~\eqref{eq:gl-dynamics}, получаем нелинейное уравнение
Шрёдингера типа Ginzburg--Landau:
\begin{equation}
\label{eq:nlse-gl}
  i\planckHbar\,\frac{\partial\Phi}{\partial t}
  = \alpha\,\Phi
  - \frac{\planckHbar^2}{2m}\,\nabla^2\Phi
  + \beta\,|\Phi|^2\Phi.
\end{equation}
В координатах плотности $|\Phi|^2$, при $\alpha=0$ на выбранном фоне, оно эквивалентно записи
\begin{equation}
\label{eq:nlse-gl-real}
  \frac{\partial\Phi}{\partial t}
  = \frac{i\planckHbar}{2m}\,\nabla^2\Phi
  - i\,\beta\,|\Phi|^2\Phi.
\end{equation}

\paragraph{Диссипация осреднения.}
Член $\nu_{\mathrm{CA}}\nabla^2\Phi$ \emph{не} следует из~\eqref{eq:gl-free-energy}:
он появляется только после учёта необратимости~$\mathcal{B}$ на~$\Tlayer$
(\cref{lem:236}, \cref{prop:237}) и дискретного coarse-graining (\cref{sec:t-cr}).
$\mathcal{F}_{\mathrm{GL}}$ задаёт \emph{обратимую} часть; $\nu_{\mathrm{CA}}$ --- кинематическая вязкость
макроуровня (ср.~диффузионное дополнение к волновому уравнению в курсе полей).
Полное уравнение на~$\Tlayer$:
\[
  \frac{\partial\Phi}{\partial t}
  = \frac{i\planckHbar}{2m}\,\nabla^2\Phi
  - i\,\beta\,|\Phi|^2\Phi
  + \nu_{\mathrm{CA}}\,\nabla^2\Phi,
\]
новый оператор на~$M$ при этом не вводится.

\section{Туннелирование и правило Бorn}

Туннелирование на~$T$ --- макроскопическое осреднение утечки через барьер.
Правило Бorn $|\Psi|^2$ --- статистика макрофильтра (\cref{prop:determinism}).
